\documentclass[aps,prd,twocolumn,10pt,nofootinbib,superscriptaddress]{revtex4-2}
\usepackage{amsmath,amssymb,graphicx}
\usepackage{booktabs}
\usepackage[colorlinks=true,linkcolor=blue,citecolor=blue,urlcolor=blue]{hyperref}

\newcommand{\FF}{\mathcal{F}}
\newcommand{\YY}{\mathcal{Y}}
\newcommand{\QQ}{\mathcal{Q}}
\newcommand{\Gt}{\tilde{G}}
\newcommand{\Msun}{M_\odot}

\begin{document}

\title{Three requirements on a relativistic home for MOND phenomenology}

\author{Carl P.\ Zimmerman}
\affiliation{Briar Creek Tech}
\date{18 August 2026 (v3)}

\begin{abstract}
Aether--scalar--tensor gravity (AeST) [Skordis \& Z\l o\'snik, Phys.\ Rev.\ Lett.\ \textbf{127},
161302 (2021)] is the only relativistic MOND-class theory known to reproduce the cosmic microwave
background power spectrum, and it contains a single free function $\FF(\YY,\QQ)$ that its authors
leave unspecified. We show that specifying it is obstructed, and we extract from the obstruction
three necessary conditions on \emph{any} relativistic completion of MOND phenomenology. Each is
established by an explicit construction that satisfies the previous conditions and fails the next.
\textbf{(R1)} The free function must depend on the gradient of the \emph{total} potential, not on
the scalar's own. With $\FF(\YY)$ the quasi-static reduction forces the anomalous acceleration
$U(y)$ to be monotone increasing, hence to saturate at some $U\to s$, hence to produce a
\emph{constant} sunward anomaly $s\,a_0$ at every planetary distance. Perihelion precession then
requires $s\le1.27\times10^{-5}$ while the radial acceleration relation requires $s\ge0.435$ ---
incompatible by $1.2$--$3.4\times10^{4}$. There is no external-field relief: computed for this
family rather than imported, it is $1.000000\times$. \textbf{(R2)} Eating
the total gradient must not destabilise the aether's vector sector. Replacing $\YY$ by
$Z\equiv J^\mu J_\mu$, the squared aether acceleration, satisfies R1 exactly --- the quasi-static
reduction becomes AQUAL, the exponential interpolation becomes legal, $\gamma_{\rm PPN}=1$ and
$G_N$ stay unchanged, the 1~AU anomaly falls to $10^{-3458.7}\,\mathrm{m\,s^{-2}}$, the FRW
background is \emph{exactly} undisturbed, and the derived $a_0(z)$ law survives bit-for-bit. It
fails on the vector modes, and it fails worst exactly where it was invoked: the vector kinetic
coefficient is $C_V=K_B-(2-K_B)J_Z$, and the same screening that voids the gap drives
$J_Z\to1$, giving $C_V=-1.50$ and $c_V^2=-0.167$ at $K_B=0.25$ --- a ghost with a gradient
instability whose growth rate $0.408\,ck$ is unbounded in the ultraviolet. The unstable region
covers the entire solar system and Oort cloud ($r<5.16\times10^4$~AU) and the entire optical body
of a galaxy ($r<79$~kpc), $e$-folding in 624~yr even at the longest wavelength for which the
reduction is controlled. \textbf{(R3)} And the repair may not be bought with a split between the
cosmological and local gravitational constants. Rescaling, $\mathcal{J}\to s\mathcal{J}$, cures
the ghost at every acceleration for $s<K_B/(2-K_B)$ --- and forces $\Gt/G_N\le0.121$, giving a primordial helium
abundance $Y_p=0.086$ against an observed $0.2449\pm0.0040$ ($-42\sigma$) and a gravitational
radiation density \emph{below that of the photons alone}. Gravitational waves do not pay this
price: the standard-siren amplitude's sensitivity is exactly zero. We also record a structural identity: AeST's action
already contains an explicit $+(2-K_B)a^\mu a_\mu$, so its apparent $c_4=0$ is an artifact of
variables; and $c_T=1$ holds exactly, for any free function.
\end{abstract}

\maketitle

\section{Setup}

AeST's action, verified against the authors' source \cite{SZ2021,SZ2022}, is
\begin{align}
S=\int d^4x\,\frac{\sqrt{-g}}{16\pi\Gt}\Big[&R-2\Lambda-\tfrac{K_B}{2}F^{\mu\nu}F_{\mu\nu}
\nonumber\\
&+2(2-K_B)J^\mu\nabla_\mu\varphi-(2-K_B)\YY \nonumber\\
&-\FF(\YY,\QQ)-\lambda(A^\mu A_\mu+1)\Big]+S_{\rm m}[g],
\label{eq:action}
\end{align}
with $\QQ=A^\mu\nabla_\mu\varphi$, $\YY=(g^{\mu\nu}+A^\mu A^\nu)\nabla_\mu\varphi\nabla_\nu\varphi$,
$F_{\mu\nu}=2\nabla_{[\mu}A_{\nu]}$, and $J^\mu=A^\nu\nabla_\nu A^\mu$ the aether's acceleration.
Matter couples to $g_{\mu\nu}$ alone.

Throughout we use the acceleration scale
\begin{equation}
a_0=\kappa c\sqrt{G\rho_\Lambda}=9.3619\times10^{-11}\,\mathrm{m\,s^{-2}},
\label{eq:a0}
\end{equation}
with an alternative footing giving $1.1279\times10^{-10}$; both are carried throughout. The
coefficient $\kappa=\tfrac12$ is \emph{fitted, not derived}: three determinations give
$0.551\pm0.043$, $0.465\pm0.076$ and $0.537\pm0.071$, combining to $0.529\pm0.034$. Write
$y=g_{\rm bar}/a_0$ and $U(y)=u/a_0$ for the anomalous acceleration.

\section{R1: the free function must eat the total gradient}

With $\FF(\YY)$, the quasi-static gradient sector diagonalises: $\Psi=\Psi_N+\chi$ with
$\nabla^2\Psi_N=4\pi\hat G\rho$, and in spherical symmetry Gauss's theorem gives the local
algebraic law $u\,J_Y(u^2)=g_{\rm bar}$ with $u=|\nabla\chi|$. The scalar's own gradient
\emph{is} the anomalous acceleration. Reconstructing the free function from a chosen interpolation
gives $J_Y(Y)=y(U)/U$, invertible iff $U$ is strictly increasing; equivalently the longitudinal
eigenvalue of $M_{ij}=J_Y\delta_{ij}+2J_{YY}u_iu_j$ is positive under the same condition. A
non-monotone $U$ means both a multi-valued $\FF$ and a longitudinal gradient ghost.

\emph{This is specific to AeST.} In AQUAL \cite{BM1984} the free function depends on the total
potential's gradient and stability requires only $dg_{\rm obs}/dg_{\rm bar}>0$ --- satisfied by the
exponential interpolation $\nu=1/(1-e^{-\sqrt y})$ \cite{MS2008}, minimum $0.968$. The same kernel
is legal in one theory and fatal in the other, purely from which gradient the free function eats.

Monotonicity plus $U/y\to0$ forces saturation, $U\to s$: a constant sunward anomaly $s\,a_0$ at
every planetary distance. By the Gauss planetary equations a constant radial perturbation gives
$\dot\varpi=s\,a_0\sqrt{1-e^2}/(na)$; against anomalous-precession limits \cite{Pitjeva,Fienga} the
worst planet requires
\begin{equation}
s\le1.27\times10^{-5}\ \text{(canonical)},\quad 1.05\times10^{-5}\ \text{(alt)}.
\end{equation}
Applying $U(2)\ge0.4$ to the family $J_Y=v/(1-v/s)$, $v=\sqrt{Y}/a_0$, gives $s\ge0.4348$; relaxing
to a global fit (rms $\le0.15$ dex on 3389 SPARC points \cite{SPARC}, $\Upsilon$ in the Spitzer
band) gives $0.157$--$0.294$ depending on treatment. Hence
\begin{equation}
\boxed{\text{GAP: }1.2\text{--}3.4\times10^{4}.}
\end{equation}

Every candidate relief was computed and none helps. \emph{The external-field effect contributes
nothing}: derived for this family by two independent routes (an $\ell=1$ penetration ODE and a flux
bound requiring no perturbation theory), the relief on the saturated anomaly is $1.000000\times$
--- the external field screens \emph{itself}, not the anomaly, and somewhere on the 1~AU sphere the
anomaly is at least $s\,a_0(1-4\times10^{-9})$. A density-dependent $a_0$ \emph{hurts}: with
suppression $f$ the constrained product is $sf$, and $f\,U_s(2/f)=0.4$ gives $0.435$ at $f=1$ but
$2.00$ at $f=0.1$, unsatisfiable below $f=0.080$ because $U\le\sqrt y$ caps the family. Solving the
non-spherical problem for a Miyamoto--Nagai disc raises the required floor by $1.07$--$1.10$.
Refitting $\Upsilon$ buys $1.34\times$.

\section{R2: eating the total gradient must not destabilise the vector sector}

R1 says: use the total gradient. AeST already contains the object that supplies it. With
$D^\mu=q^{\mu\nu}\nabla_\nu\varphi$ one has the identity, verified three times independently and
requiring no computer algebra,
\begin{align}
&2(2-K_B)J\!\cdot\!\nabla\varphi-(2-K_B)\YY \nonumber\\
&\qquad\equiv +(2-K_B)Z-(2-K_B)|D-J|^2,
\label{eq:identity}
\end{align}
with $Z\equiv J^\mu J_\mu$. \emph{AeST therefore already contains an explicit
$+(2-K_B)a^\mu a_\mu$: its apparent $c_4=0$ is an artifact of variables.} Moreover
$J^i=\nabla^i\Psi$ with coefficient exactly unity, with $Z$'s leading coefficient independent of
$\varphi$, $\Phi$ and the aether perturbation --- so $Z$ \emph{is} the total potential gradient
squared.

Replacing $\FF(\YY,\QQ)$ by $\FF(Z,\QQ)$ then works exactly as R1 requires. The aether-longitudinal
equation becomes $v=\nabla\Psi$ pointwise, the $(1+J_Y)$ factor disappears, and the reduction is
\emph{AQUAL exactly}, with $\mu(g_{\rm obs})=J_Z(g_{\rm obs}^2)$ and $\hat G$ unchanged. The
exponential kernel becomes legal (min $dx/dy=0.968$, against min $dU/dy=-0.032$ under $\FF(\YY)$
for the same kernel), and the 1~AU anomaly becomes $10^{-3458.7}\,\mathrm{m\,s^{-2}}$ --- some
$10^{3445}$ below the Sereno--Jetzer bound \cite{SerenoJetzer}, replacing an $s\,a_0$ that sat
$1279\times$ above it. Further, $c_T^2=1$ exactly for any free function, and $G_N$ is finite and
equal to the unmodified value.

This construction does everything R1 asks, and several things it does not. The FRW background
is \emph{exactly} undisturbed: the comoving aether congruence is geodesic, so $J^\mu=0$
identically for arbitrary lapse, hence $\bar Z=0$ and $\FF(0,\QQ)=K(\QQ)$ --- the same background
function AeST's own $\FF(\YY,\QQ)$ reduces to --- and the new term's stress contribution carries at
least one factor of $J$ and vanishes too. Friedmann, $w=-1$, the dust mode, $Q(a)$ and therefore
the derived $a_0(z)$ law are bit-for-bit the AeST objects. In the static sector one finds AQUAL
with $\mu=J_Z$, $\gamma_{\rm PPN}=1$ and $G_N=\hat G$ finite. At linear order
$J_i=\nabla_i E$ exactly, with $E=\dot\alpha+\Psi$ the theory's own mixing variable, so
$Z=|\nabla E|^2/a^2$ is second order; and since $F^{\mu\nu}F_{\mu\nu}=-2Z$ at $O(\delta^2)$, any
analytic linear-in-$Z$ piece of $\FF$ is a mere relabelling $K_B\to K_B-\FF_Z(0)$, with
$\FF_Z(0)=0$ on the determined family because every interpolation's deep-MOND end has
$\mu_{\rm phys}\to0$.

Its cosmological cost is real but payable, and we state it because it is the one place the
framework's own machinery does unexpected work. Eating the \emph{total} gradient switches the MOND
branch on at recombination, where a CLASS run puts the acoustic-band gravitational acceleration at
only $y=7.6$--$11.9$ absent running; the resulting $\epsilon=1-\mu(y)\simeq1/2y$ is a few per cent
kernel-robustly, worth $18$--$112\sigma$ aggregated over $\ell=2$--$2500$ at unit response. The CMB
tolerates it only if $a_0({\rm rec})/a_0(0)\le9.7\times10^{-3}$ to $3.7\times10^{-1}$ across eight
forks --- and the suppression this framework independently derives, $0.0060$, clears the tightest
fork and the loosest, neither number having been tuned to the other. (Delivering it needs
$\nu_0\ge5.6\times10^{-9}$ to $8.2\times10^{-6}$, which exceeds the independent RAR ceiling
$2.36\times10^{-6}$ on two of eight forks --- exactly the $\alpha=1$ kernel under the conservative
reading --- while the exponential kernel this construction re-legalises clears on every reading.
The interlock scales as $S^2$ in the uncomputed CMB response $S=|d\ln C_\ell/d\epsilon|$.)

\emph{It dies on the vector modes, and it dies worst exactly where it was invoked.} The
transverse-vector kinetic coefficient is
\begin{equation}
C_V=K_B-(2-K_B)J_Z,\qquad c_V^2=K_B/C_V,
\label{eq:CV}
\end{equation}
so the theory is ghost-free only for $J_Z<K_B/(2-K_B)$. But $J_Z=\mu_{\rm phys}$ is precisely the
interpolation function, and the Newtonian screening that deletes the ephemeris gap is the
statement $\mu_{\rm phys}\to1$. At $K_B=0.25$ this gives $C_V=-1.50$ and $c_V^2=-0.167$: a ghost
carrying a gradient instability, growth rate $0.408\,ck$, unbounded in the ultraviolet. Converting
$J_Z<K_B/(2-K_B)$ into a radius, the unstable region covers the entire solar system and Oort cloud
($r<5.16\times10^{4}$~AU canonical, $4.70\times10^4$ alt) and the entire optical body of a
$10^{11}M_\odot$ galaxy ($r<79.2$~kpc canonical, $72.1$ alt). Even at the longest wavelength for
which the reduction is controlled, $\lambda=c^2/g$ at 1~AU, the $e$-folding time is 624~yr. No
cutoff choice saves it, and only $K_B\ge1$ removes the bound at all.

The structure of the failure is the point: \emph{the same factor $\mu_{\rm phys}\to1$ that makes
the solar-system screening work is what makes $C_V$ most negative}. Escaping R1 and destabilising
the vector sector are one operation, not two. AeST itself, with $\FF(\YY)$ and no $Z$-term, has
$C_V=K_B>0$ and is healthy here; the liability belongs to this construction alone.

\section{R3: and the repair may not be bought with a $\Gt/G_N$ split}

Equation~\eqref{eq:CV} names its own repair. Both sector conditions constrain only the total
$Z$-coefficient: static attraction needs it positive, vector no-ghost needs it below $K_B$.
Rescaling the free function, $J\to sJ$, satisfies both at every acceleration for
$s<K_B/(2-K_B)$; the window is $0<s<0.138$ at $K_B=0.25$, and the solar-system screening is
untouched, being $s$-independent.

The price is that with $\FF(Z)$ the non-$\FF$ static Lagrangian is annihilated identically by
$v=\nabla\Psi$, making the free function's normalisation the sole source of Newton's constant:
$G_N=\hat G/s$. The no-ghost condition is therefore equivalent to $2\Gt<K_B G_N$, i.e.
\begin{equation}
\Gt/G_N\le0.121\ (K_B=0.25),\qquad 0.048\ (K_B=0.10),
\end{equation}
against AeST's $0.875$ --- an $8.3\times$ to $20.7\times$ split. It cannot be paid.

\emph{Big-bang nucleosynthesis.} $Y_p=0.086$ at the no-ghost ceiling against an observed
$0.2449\pm0.0040$: $-42\sigma$ from the same calculation's $\Gt/G_N=1$ control. A closure-free
bound assuming no nucleosynthesis model at all still gives $Y_p\le0.176$, i.e.\ $-17\sigma$.

\emph{The CMB, and this is an inconsistency rather than a tension.} The gravitational sector's
radiation density comes out at $0.2046$ of the thermal photon density --- \emph{below the photons
alone}, whose value FIRAS \cite{FIRAS} fixes with no gravity anywhere in the chain. The equivalent
$N_{\rm eff}=-3.503$. The hard floor is missed by $4.9\times$ ($K_B=0.25$) to $12.2\times$
($K_B=0.10$), and the ceiling sits $139\sigma$ from viability in cosmic-variance units.

\emph{Gravitational waves do not pay this price, and we state that against our own interest.} The
tensor sector is exactly Einstein--Hilbert: on FRW plus a transverse-traceless (TT) mode,
$J^\mu=0$, $Z=0$, $F_{\mu\nu}=0$ and $\YY=0$ all hold identically and to all orders in $h$, so the
dark sector carries no anisotropic stress for any free function, giving $c_T=1$ and $\alpha_M=0$
exactly. The standard-siren amplitude's sensitivity to $s$ is then exactly zero: $\Gt$ enters the
strain and $G_N$ the orbit, and they cancel in $d_L^{\rm GW}$. GW170817 \cite{GW170817} places
\emph{no bound} on $s$. Two liabilities the same derivation creates are recorded but \emph{not
scored}: the inferred chirp mass rescales as $r^{3/5}$ ($4.13\,\Msun$ for GW170817 at the ghost
edge, which needs the theory's own strong-field stellar structure to close), and the tensor-channel
orbital decay is $\dot P_b/\dot P_b^{\rm GR}=r$, an $8\times$ deficit against B1913+16
\cite{WeisbergHuang}. We decline to score the second because the identical ledger applied to
\emph{unmodified} AeST would exclude it at $77\sigma$ --- the tensor-only flux ledger is evidently
incomplete, and the total radiated flux is not computed here. The usable result is a constraint on
the shape of any repair: the missing channels must supply the same $1-r$ fraction at every orbital
period, since $\dot P_b/\dot P_b^{\rm GR}$ is measured to $0.16\%$ in a 7.75~h orbit and to
$1.3\times10^{-4}$ in a 2.45~h one \cite{Kramer2021}, and a dipole enters at relative order
$(c/v)^2$ and cannot be period-independent.

\section{Testable predictions}

A no-go is tested by construction rather than by observation, but the derivations above carry four
consequences that data can decide.

\textbf{P1: $\alpha_M=0$ and $c_T=1$, exactly, for any free function.} This is the strongest
prediction here and the one least dependent on anything we have chosen. On FRW plus a TT mode,
$J^\mu=0$, $Z=0$, $F_{\mu\nu}=0$ and $\YY=0$ hold identically and \emph{to all orders in} $h$, so
every non-Einstein invariant is $h$-independent and the dark sector's TT stress is
$T^x{}_x-T^y{}_y=0$. The tensor sector is therefore exactly Einstein--Hilbert with
$M_*^2=1/8\pi\Gt$ constant. AeST and every construction in this paper are consequently
indistinguishable from general relativity in gravitational-wave propagation --- no running Planck
mass, no modified friction, no $c_T$ deviation, at any redshift and for any $K_B$, $s$ or $a_0$
footing. GW170817's $|c_T-1|\lesssim10^{-15}$ is passed trivially rather than fitted. \emph{A
confirmed $\alpha_M\neq0$ from LISA or Einstein Telescope standard sirens falsifies this entire
class in one measurement}, and cannot be absorbed by any choice of free function.

\textbf{P2: the orbital-decay deficit must be period-independent.} If the tensor-only flux ledger
of Sec.~IV is incomplete, as we argue, the missing channels must supply the same $1-r$ fraction of
the flux at every orbital period. A dipole channel enters at relative order $(c/v)^2$ and therefore
cannot be period-independent; the repair must renormalise the quadrupole coefficient itself. Any
future binary showing a period-\emph{dependent} deviation from the GR quadrupole formula excludes
the repair, and with it the aether-vector rescaling.

\textbf{P3: the $\FF(\YY)$ class predicts a constant sunward anomaly with a fixed planet ordering,}
$\dot\varpi=s\,a_0\sqrt{1-e^2}/(na)$, falling with semi-major axis in a way no fitted $\dot\varpi$
per planet can mimic, since a single $s$ sets all of them. This is already excluded (Sec.~II),
which is \emph{why} R1 exists; a future ephemeris solution reporting a coherent residual of this
shape would revive rather than kill it.

\textbf{P4: the $Z$-form predicts no solar-system deviation at all,} at any measurable level:
the 1~AU anomaly is $10^{-3458.7}\,\mathrm{m\,s^{-2}}$ and is $s$-independent. Any claimed
solar-system MOND detection excludes it outright.

These are separate from framework-level predictions that this paper leaves untouched, since they
belong to the normalisation \eqref{eq:a0} rather than to any completion: a hash-frozen wide-binary
velocity ratio decided by Gaia DR4; a declining $a_0(z)$, which puts the MOND regime \emph{off} at
recombination at $0.0060$ of its present value and is falsified by any robust $a_0$ evolution below
$z\sim5$; and $\nu_0\le2.36\times10^{-6}$, i.e.\ no environmental variation of $a_0$ at a level
current data could see.

\section{What is and is not claimed}

\emph{Claimed.} The three requirements, each with an explicit construction that satisfies its
predecessors and fails it. The identity \eqref{eq:identity} and its corollary that AeST's $c_4=0$
is a choice of variables --- and, more sharply, that $c_4^{\rm eff}$ is \emph{sector-dependent},
so ``$c_{14}=2$'' is not a well-defined statement about this theory. That $c_T=1$ holds exactly in
AeST and in the $\FF(Z)$ theory for any free function. That the monotonicity obstruction of R1 is
genuinely escapable, and what escaping it costs. As a subsidiary but independent result: that the
$Z$-form's cosmological cost is met by the framework's own derived $a_0(z)$ suppression on all
eight forks --- which does not save it, R2 being fatal, but does establish that the CMB was never
the binding constraint on this construction.

\emph{Not claimed.} That the requirements are sufficient --- they are necessary conditions
extracted from three failures, and a completion satisfying all three is not thereby viable. That
the $Z$-form's CMB verdict is settled: it is conditional on an uncomputed response and on two of
eight forks, and is in any case not what kills it. That the scalar dynamical modes, $\alpha_1$ or
$\alpha_2$ have been computed for the $Z$-form; they have not. That
$\kappa=\tfrac12$ is derived; it is fitted, and must be quoted with its $H_0$, the distance-free
determination carrying an unpriced $\sim7\%$ systematic. That dark matter is absent: $\Omega_{\rm
dm}$ is full here and the claim is ``no dark-matter \emph{particle}''. Whether the dark sector's
dust remains bound inside galaxies is unresolved and untouched by any of this --- it is a
$\QQ$-sector problem, and all three constructions modify the $\YY$-sector.

\emph{Unaffected by any of the above}, being independent of which completion carries them: the
normalisation \eqref{eq:a0}, the radial acceleration relation at 0.108 dex on 175 SPARC galaxies,
weak lensing from 40 kpc to 2.2 Mpc with no dark component, the baryonic Tully--Fisher relation,
and a hash-frozen wide-binary prediction decided by Gaia DR4.

\emph{What a viable completion must supply.} A screening mechanism that is not a non-monotone
interpolation (R1); one whose approach to the Newtonian limit does not simultaneously drive a
kinetic coefficient negative (R2); and one that does not separate the cosmological and local
gravitational constants (R3). R2 is the sharpest of the three as a design constraint, because it
is not a statement about any particular invariant: it says that whatever object the free function
eats, the \emph{same limit} that produces Newtonian screening must not be the limit that
destabilises another sector. Superfluid constructions \cite{BK2015}, in which the MOND force
exists only inside a condensate phase and the Newtonian exterior is a different phase rather than
a limit of the same function, evade that coupling by construction and are the obvious next
candidate.

\emph{A note on the process.} Version~1 of this paper attributed the $Z$-form's death to a
cosmological degeneracy $\FF_Z(0)=K_B$; version~2, correcting that, over-swung and reported the
construction as viable. Both were wrong. The first misread ``degenerate with the $K_B$ term''
(indistinguishable from it) as ``the equation degenerates''; the second, having removed a false
death, failed to check that a true one was already on record in the vector sector. The correct
statement is the one above: the cosmology is payable, and the vector ghost is fatal. Both errors
are recorded, dated, with their directions stated, in \texttt{RETRACTIONS.md}.

\begin{acknowledgments}
Every number in this paper is produced by a committed self-checking script at
\url{https://github.com/carlzimmerman/zimmerman-formula}; each prints numbered checks and exits
non-zero on failure. Withdrawn claims are recorded, dated, in \texttt{RETRACTIONS.md} --- including
one in this paper's own drafting, where a gravitational-wave exclusion was written from an
unverified number and corrected against the committed script before filing. The contribution here
is the three requirements, the identity \eqref{eq:identity}, and the constructions that establish
them --- not the underlying theory, which is Skordis and Z\l o\'snik's.
\end{acknowledgments}

\begin{thebibliography}{99}
\bibitem{SZ2021} C.~Skordis and T.~Z\l o\'snik, Phys.\ Rev.\ Lett.\ \textbf{127}, 161302 (2021),
  arXiv:2007.00082.
\bibitem{SZ2022} C.~Skordis and T.~Z\l o\'snik, arXiv:2109.13287.
\bibitem{BM1984} J.~Bekenstein and M.~Milgrom, Astrophys.\ J.\ \textbf{286}, 7 (1984).
\bibitem{MS2008} M.~Milgrom and R.~H.~Sanders, Astrophys.\ J.\ \textbf{678}, 131 (2008),
  Eq.~(13) at $\alpha=\tfrac12$.
\bibitem{Pitjeva} E.~V.~Pitjeva and N.~P.~Pitjev, Mon.\ Not.\ R.\ Astron.\ Soc.\ \textbf{432},
  3431 (2013).
\bibitem{Fienga} A.~Fienga \textit{et al.}, Celest.\ Mech.\ Dyn.\ Astron.\ \textbf{111}, 363 (2011).
\bibitem{SPARC} F.~Lelli, S.~S.~McGaugh, and J.~M.~Schombert, Astron.\ J.\ \textbf{152}, 157 (2016).
\bibitem{SerenoJetzer} M.~Sereno and Ph.~Jetzer, Mon.\ Not.\ R.\ Astron.\ Soc.\ \textbf{371}, 626
  (2006).
\bibitem{FIRAS} D.~J.~Fixsen \textit{et al.}, Astrophys.\ J.\ \textbf{473}, 576 (1996).
\bibitem{GW170817} B.~P.~Abbott \textit{et al.} (LIGO/Virgo), Phys.\ Rev.\ Lett.\ \textbf{119},
  161101 (2017).
\bibitem{WeisbergHuang} J.~M.~Weisberg and Y.~Huang, Astrophys.\ J.\ \textbf{829}, 55 (2016).
\bibitem{Kramer2021} M.~Kramer \textit{et al.}, Phys.\ Rev.\ X \textbf{11}, 041050 (2021).
\bibitem{BK2015} L.~Berezhiani and J.~Khoury, Phys.\ Rev.\ D \textbf{92}, 103510 (2015).
\bibitem{opt1cmb} Reproduced by \texttt{real\_research/reviews/opt1\_cmb\_2026.py}
  (36/36 checks), which carries the CLASS run, the eight-fork table and the $S$-threshold.
\bibitem{FJ2006} B.~Z.~Foster and T.~Jacobson, Phys.\ Rev.\ D \textbf{73}, 064015 (2006).
\end{thebibliography}

\end{document}
